\hypertarget{index_intro_sec}{}\section{Introduction}\label{index_intro_sec}
Welcome to a brief introduction of a project executed by a student of the university of Kassel. ~\newline
 The goal of this project is to create reliable software that controlls the robot to drive three laps on a track (\href{track_img.jpg}{\tt {\bfseries example track}}) consisting of a black line, autonomously. ~\newline
There is a black square placed on the track which acts as the startfield. ~\newline
 The robot has three sensors which enable it to detect the black line aswell as the startfield so it can take action accordingly. ~\newline
~\newline
In addition to that the robot starts driving only after being placed on the startfield which initiates a 15s countdown. ~\newline
Upon being moved the countdown shall be interrupted and the robot has to look for the startfield again. ~\newline
 The countdown has the robot blinking three leds with 5hz which are placed on top of the roboter. ~\newline
When the robots completes 3 laps it stops and a 60s countdown starts, analog to the first countdown.\hypertarget{index_install_sec}{}\section{Installation}\label{index_install_sec}
Installation and handelling is not as compilcated as one might think!
\begin{DoxyItemize}
\item First, one has to create a U\+SB connection between the computer and the robot.
\item Open the terminal and navigate to the file with the code using \char`\"{}cd\char`\"{} (= change directory), or open the terminal in the desired file folder by clicking on the right mouse pad and choosing \char`\"{}\+Open Terminal Here\char`\"{}.
\item Depending on which option you wish, chose one of the following compiler options\+:
\end{DoxyItemize}

The \href{transition_map.pdf}{\tt {\bfseries Transition Map}}.\hypertarget{index_Hardware}{}\section{Hardware}\label{index_Hardware}
\hypertarget{index_Table_Of_Contents}{}\section{Table\+\_\+\+Of\+\_\+\+Contents}\label{index_Table_Of_Contents}
\hypertarget{index_States}{}\subsection{descrption}\label{index_States}
to implement new states\hypertarget{index_Calibration}{}\section{Calibration}\label{index_Calibration}
\hypertarget{index_compiler_options}{}\section{Compiler Options}\label{index_compiler_options}

\begin{DoxyItemize}
\item {\bfseries make}\+: compiles, flashes and cleans up files that are no longer useful. Debug modus is disabled.
\item {\bfseries make debug}\+: debug modus is enabled. ~\newline

\item {\bfseries make remote}\+: This modus enabled the remote control of the robot via cutecom. Debug modus is disabled. ~\newline

\item {\bfseries make documentation}\+: This modus only updates the doxygen documentation. . ~\newline
 {\bfseries Attention} if you received the file, usually a doxyfile should already exist. However, if there is no html or latex file in the file you received, first type \textquotesingle{}doxygen -\/g doxyfile\textquotesingle{} in the terminal. After that open the doxyfile in nano (\textquotesingle{}nano doxyfile\textquotesingle{} in the terminal) and search for \textquotesingle{}O\+P\+T\+I\+M\+I\+Z\+E\+\_\+\+O\+U\+T\+P\+U\+T\+\_\+\+F\+O\+R\+\_\+C\textquotesingle{} (use the explanation at the bottom of the screen in nano for short cuts). Now you have to type \textquotesingle{}Y\+ES\textquotesingle{} behind the \textquotesingle{}=\textquotesingle{}-\/symbol and save the doxyfile.
\end{DoxyItemize}\hypertarget{index_Amazing_features}{}\section{Extra features}\label{index_Amazing_features}
The robot not only drives three rounds on a circular track, has two reliable and exact countdowns using different timer and prescaler to show the robot\textquotesingle{}s (and coder\textquotesingle{}s \+:) ) full potential, but also {\bfseries only} uses timer instead of delays and has a working remote control! ~\newline
To use the remote control, you first have to enable it with the compile option \textquotesingle{}make remote\textquotesingle{}. ~\newline
 Open cutecom and create a bluetooth connection. ~\newline
 Then type a direction with the desired time in ms in the text box and send it. ~\newline
 The robot will follow your order as long as the time you sent is not expired or you have not send a new order. Here is a list of the directions available\+: ~\newline

\begin{DoxyItemize}
\item F\+: forwards
\item B \+: backwards
\item l \+: soft left
\item r \+: soft right
\item L \+: hard right
\item R \+: hard left
\item Q \+: do not move ~\newline
 For example, if you type and send \textquotesingle{}F100\textquotesingle{} the robot will move forwards for 100 ms.
\end{DoxyItemize}\hypertarget{index_About_me}{}\section{About the author}\label{index_About_me}
To tell you a little bit more about myself\+: My name is Klara M. Gutekunst, i am currently at the end of my second semester studying computer science at the university of Kassel. ~\newline
This semester we had the chance to choose a project in a variety of hands-\/on coding training, including programming a robot. As I have never worked with code before I started studying, ~\newline
 I wanted to bring what I have learned so far to good use and actually see what some lines of code can do in real life. I have been not disappointed, as you can see\+: ~\newline
The work really paid off and at the end my robot slowely but confidently moves around on the track!\hypertarget{index_further_information}{}\section{Optionial information}\label{index_further_information}
If you wish to read more about the code and why certain code is used the way it is, just explore the option given above in the main page. ~\newline
If you have any suggestions for improvement, do not hesitate and contact me\+: \href{mailto:klara.gutekunst.kg@gmail.com}{\tt klara.\+gutekunst.\+kg@gmail.\+com} ~\newline
I hope you enjoy this project as much as I did\hypertarget{Mainpage.md_Roadmap}{}\section{Roadmap}\label{Mainpage.md_Roadmap}
